% =========================================================
% Stage 117 S117-A (Hypothetical/Test Unit)
% =========================================================
\part{Stage 117 --- Hypothetical Advanced Structures}
\chapter{S117-A --- Rollback Union-Find (DSU with Undo)}

\begin{notebox}[가상 Learning Unit]
S117-A는 현재 실제 커리큘럼에 존재하지 않는 가상 Unit이다. 이 자료는 학습자료 운영 구조와 기록 절차를 시험하기 위해 만든다. 따라서 이 Unit의 진행은 실제 S1-A 숙달도, Review Debt, 다음 진도에 영향을 주지 않는다. 다만 여기에서 설명하는 Rollback Union-Find 자체는 실제 competitive programming에서 사용하는 기법이다.
\end{notebox}

% =========================================================
% Part A
% =========================================================
\section{Part A --- 위치와 선수지식}

\subsection{A1. 가상의 현재 위치}
이 Unit은 다음 내용을 이미 학습했다고 가정한다.
\begin{itemize}
\item 기본 Union-Find / Disjoint Set Union (DSU)
\item \code{parent} 배열과 root 개념
\item union by size 또는 union by rank
\item stack 또는 \code{vector}를 이용한 변경 이력 저장
\item $O(\log N)$, amortized complexity의 기본 의미
\end{itemize}

일반 DSU는 원소들을 여러 disjoint set으로 관리하면서 다음 연산을 빠르게 수행한다.
\begin{itemize}
\item 두 원소가 같은 집합에 속하는지 확인
\item 서로 다른 두 집합을 합치기
\end{itemize}

S117-A에서는 여기에 새로운 요구를 추가한다.

\begin{quote}
``방금까지 수행한 union들을 특정 과거 시점까지 되돌릴 수 있어야 한다.''
\end{quote}

\subsection{A2. 왜 별도의 기법이 필요한가}
일반 DSU는 과거 상태를 복구하는 기능을 기본적으로 제공하지 않는다. 특히 path compression은 \code{find()}를 호출하는 것만으로도 여러 \code{parent} 값을 바꿀 수 있다. 따라서 단순히 ``마지막 union 하나만 기억하면 되겠지''라고 생각하면 실제 변경량을 놓칠 수 있다.

Rollback DSU의 핵심 설계는 반대 방향이다.
\begin{itemize}
\item path compression을 사용하지 않는다.
\item union by size/rank로 tree 높이를 제한한다.
\item 성공한 union이 변경한 소수의 값만 history에 저장한다.
\item history를 역순으로 되돌려 과거 상태를 복구한다.
\end{itemize}

\subsection{A3. 이후 연결되는 주제}
Rollback DSU는 다음과 같은 고급 문제의 기반이 될 수 있다.
\begin{itemize}
\item offline dynamic connectivity
\item 시간 구간별 edge 활성화와 segment tree over time
\item divide-and-conquer over queries
\item 탐색 과정에서 상태를 변경했다가 재귀 복귀 시 복원하는 문제
\end{itemize}

\subsection{A4. Core Decision Boundaries}
이 Unit은 여러 구현 선택을 구분해야 하는 multi-tool 성격을 갖는다.
\begin{enumerate}
\item \textbf{일반 DSU vs Rollback DSU}: 과거 상태 복원이 필요한가?
\item \textbf{path compression 사용 vs 미사용}: rollback이 필요한 상태에서 \code{find()}가 광범위한 변경을 일으켜도 되는가?
\item \textbf{전체 상태 복사 vs 변경 로그}: snapshot마다 $O(N)$ 복사가 필요한가, 실제로 바뀐 필드만 기록할 수 있는가?
\end{enumerate}

이 가상 Unit의 Immediate Coverage Floor는 실제 평가를 실시하게 될 경우 다음 두 종류의 evidence를 요구한다고 가정한다.
\begin{itemize}
\item rollback DSU의 canonical implementation 하나
\item ordinary DSU와 rollback DSU 중 무엇이 필요한지 선택해야 하는 contrastive problem 하나
\end{itemize}

% =========================================================
% Part B
% =========================================================
\section{Part B --- 학습 목표}

Part D까지 학습한 뒤에는 자료를 보지 않고 최소한 다음을 수행할 수 있어야 한다.
\begin{enumerate}
\item rollback이 필요한 문제에서 ordinary DSU의 한계를 설명할 수 있다.
\item path compression을 제거하고 union by size를 유지하는 이유를 설명할 수 있다.
\item \code{parent}, \code{size}, history를 이용한 Rollback DSU를 C++17/20으로 구현할 수 있다.
\item 현재 history 크기를 snapshot token으로 저장하고 해당 시점까지 rollback할 수 있다.
\item 성공한 union 하나를 복원하는 데 어떤 정보가 필요한지 설명할 수 있다.
\item \code{find}, \code{unite}, \code{snapshot}, \code{rollback}의 시간복잡도를 분석할 수 있다.
\item 전체 배열 복사 방식과 변경 로그 방식의 시간·공간 trade-off를 비교할 수 있다.
\end{enumerate}

% =========================================================
% Part C
% =========================================================
\section{Part C --- 개념과 원리}

\subsection{C1. 해결하려는 문제}
$N$개의 원소가 처음에는 모두 서로 다른 집합에 속한다고 하자. 다음 세 종류의 작업이 필요하다.
\begin{itemize}
\item 두 집합을 합친다.
\item 두 원소가 현재 같은 집합인지 확인한다.
\item 과거에 저장한 snapshot 시점으로 상태를 되돌린다.
\end{itemize}

문제는 현재 상태만 빠르게 유지하는 것이 아니라 \textbf{과거 상태도 효율적으로 복구}하는 것이다.

\subsection{C2. 가장 단순한 방법: snapshot마다 전체 상태 복사}
가장 직접적인 방법은 snapshot을 만들 때마다 \code{parent}와 \code{size} 배열 전체를 복사하는 것이다.

원소 수가 $N$이면 snapshot 하나당 $O(N)$ 시간과 $O(N)$ 추가 공간이 든다. snapshot이 $S$개라면 최악의 경우 $O(NS)$ 공간까지 필요할 수 있다.

작은 입력에서는 가능하지만, snapshot이 자주 발생하는 문제에서는 병목이 된다.

\subsection{C3. 핵심 관찰: successful union은 상태를 조금만 바꾼다}
path compression을 사용하지 않고 union by size만 사용한다고 하자.

서로 다른 root \code{a}, \code{b}를 합칠 때, \code{sz[a] >= sz[b]}가 되도록 정리하면 실제 변경은 다음 두 가지뿐이다.
\begin{lstlisting}
parent[b] = a;
sz[a] += sz[b];
\end{lstlisting}

즉 successful union 하나가 바꾸는 정보는 $O(1)$개다.

따라서 변경 전 상태를 history에 기록해 두면 전체 배열을 복사하지 않아도 된다.

\subsection{C4. 왜 path compression을 사용하지 않는가}
일반 DSU의 대표 최적화는 path compression이다.
\begin{lstlisting}
int find(int x) {
    if (parent[x] == x) return x;
    return parent[x] = find(parent[x]);
}
\end{lstlisting}

이 코드는 빠르지만 \code{find()} 한 번이 경로 위의 여러 \code{parent} 값을 동시에 수정한다. rollback을 위해서는 그 모든 변경을 정확히 기록해야 한다.

Rollback DSU에서는 구현과 복구 비용을 단순하게 유지하기 위해 보통 path compression을 사용하지 않는다. 대신 union by size/rank를 유지한다.

\begin{keybox}[핵심 trade-off]
Rollback DSU는 ``현재 연산을 최대한 빠르게''보다 ``변경을 적게 만들어 과거 상태를 쉽게 복원''하는 쪽을 선택한다.
\end{keybox}

\subsection{C5. union by size만으로 tree 높이가 제한되는 이유}
작은 tree의 root를 큰 tree의 root 아래에 붙인다고 하자. 어떤 노드의 depth가 1 증가하려면 그 노드가 속한 tree가 더 큰 tree와 합쳐져야 한다.

작은 쪽을 큰 쪽에 붙이므로 depth가 증가할 때마다 그 노드가 속한 집합의 크기는 최소 두 배가 된다.

따라서 한 노드의 depth는 최대 $O(\log N)$번만 증가할 수 있고, path compression이 없어도 \code{find()}는 $O(\log N)$에 수행할 수 있다.

\subsection{C6. history에 무엇을 저장하는가}
다음 successful union을 생각하자.
\begin{lstlisting}
parent[b] = a;
sz[a] += sz[b];
\end{lstlisting}

되돌리려면 최소한 다음 정보가 필요하다.
\begin{itemize}
\item child가 된 root \code{b}
\item parent가 된 root \code{a}
\item union 직전의 \code{sz[a]}
\end{itemize}

그러면 rollback은 다음처럼 수행할 수 있다.
\begin{lstlisting}
parent[b] = b;
sz[a] = old_size_a;
\end{lstlisting}

\subsection{C7. snapshot은 history의 크기다}
history가 \code{vector<Change>}라면 snapshot 시점에 배열 전체를 복사할 필요가 없다.

\begin{lstlisting}
int token = history.size();
\end{lstlisting}

이 token은 ``이 시점에는 successful union 기록이 몇 개 있었는가''를 의미한다.

나중에 다음처럼 history의 크기가 token이 될 때까지 되돌린다.
\begin{lstlisting}
while (history.size() > token) {
    undo_one_change();
}
\end{lstlisting}

\subsection{C8. 유지해야 할 invariant}
Rollback DSU에서 중요한 invariant는 다음과 같다.
\begin{itemize}
\item root $r$은 항상 \code{parent[r] == r}을 만족한다.
\item \code{sz[r]}는 root $r$이 대표하는 현재 집합의 크기다.
\item history의 각 항목은 successful union 하나가 만든 변경을 정확히 복구할 수 있다.
\item rollback은 history의 역순으로 수행한다.
\end{itemize}

마지막 조건이 특히 중요하다. 나중에 수행한 union을 먼저 취소해야 이전 union 당시의 구조가 다시 나타난다.

\subsection{C9. C++17 구현 골격}
\begin{lstlisting}
struct RollbackDSU {
    struct Change {
        int child_root;
        int parent_root;
        int parent_size_before;
    };

    vector<int> parent;
    vector<int> sz;
    vector<Change> history;

    RollbackDSU(int n) : parent(n + 1), sz(n + 1, 1) {
        iota(parent.begin(), parent.end(), 0);
    }

    int find(int x) const {
        while (parent[x] != x) {
            x = parent[x];
        }
        return x;
    }

    bool unite(int a, int b) {
        a = find(a);
        b = find(b);
        if (a == b) return false;

        if (sz[a] < sz[b]) swap(a, b);

        history.push_back({b, a, sz[a]});
        parent[b] = a;
        sz[a] += sz[b];
        return true;
    }

    int snapshot() const {
        return static_cast<int>(history.size());
    }

    void rollback(int token) {
        while (static_cast<int>(history.size()) > token) {
            Change c = history.back();
            history.pop_back();

            parent[c.child_root] = c.child_root;
            sz[c.parent_root] = c.parent_size_before;
        }
    }
};
\end{lstlisting}

여기서는 같은 집합끼리 \code{unite()}를 호출해도 상태가 변하지 않으므로 history에 아무것도 추가하지 않는다. snapshot이 history 크기 자체이기 때문에 이것으로 충분하다.

\subsection{C10. 시간복잡도와 공간복잡도}
union by size만 사용하고 path compression을 사용하지 않는 경우:
\begin{itemize}
\item \code{find(x)}: $O(\log N)$ worst-case
\item \code{unite(a,b)}: 두 번의 find가 지배하므로 $O(\log N)$
\item \code{snapshot()}: $O(1)$
\item successful union 하나를 되돌리기: $O(1)$
\item \code{rollback(token)}: 실제로 취소하는 successful union 수를 $K$라 하면 $O(K)$
\end{itemize}

초기 배열은 $O(N)$이고, successful union history를 $U$개 저장하면 총 저장공간은 $O(N+U)$다.

\subsection{C11. 언제 사용하고 언제 사용하지 않는가}
Rollback DSU가 자연스러운 경우:
\begin{itemize}
\item 탐색 중 union을 적용했다가 재귀 복귀 시 취소해야 함
\item offline query에서 특정 시간 구간에만 edge가 활성화됨
\item snapshot/rollback이 빈번하고 전체 상태 복사가 너무 비쌈
\end{itemize}

반대로 삭제나 rollback이 전혀 없고 현재 connectivity만 필요하다면 ordinary DSU가 더 단순하다. 이 경우에는 path compression까지 사용한 일반 DSU가 더 적합하다.

% =========================================================
% Part D
% =========================================================
\section{Part D --- Worked Example}

\subsection*{D1. 문제}
\begin{problembox}{Worked Example --- 임시 네트워크 연결}
$1$번부터 $N$번까지의 서버가 처음에는 모두 분리되어 있다. 다음 명령을 순서대로 처리하라.
\begin{itemize}
\item \code{LINK a b}: 두 서버가 속한 네트워크를 합친다.
\item \code{SAVE}: 현재 상태를 하나의 snapshot으로 저장한다. 예제에서는 가장 최근 SAVE만 복구 대상으로 사용한다.
\item \code{RESTORE}: 가장 최근 SAVE 시점의 상태로 되돌린다.
\item \code{SAME a b}: 현재 두 서버가 같은 네트워크에 있으면 \code{YES}, 아니면 \code{NO}를 출력한다.
\end{itemize}
\end{problembox}

예를 들어 다음 명령을 처리한다고 하자.
\begin{TextBlock}
N = 5
LINK 1 2
SAVE
LINK 2 3
SAME 1 3
RESTORE
SAME 1 3
LINK 4 5
SAME 4 5
\end{TextBlock}

기대 출력은 다음과 같다.
\begin{TextBlock}
YES
NO
YES
\end{TextBlock}

\subsection*{D2. Brute Force}
SAVE마다 \code{parent}와 \code{size} 배열 전체를 복사해 둘 수 있다. 그러면 RESTORE는 저장된 배열을 다시 복사하면 된다.

이 방식은 구현은 단순하지만 SAVE/RESTORE마다 $O(N)$ 비용이 들며 snapshot이 많으면 메모리도 크게 증가한다.

\subsection*{D3. Complexity}
전체 상태 복사 방식에서 SAVE 하나는 $O(N)$ 시간과 $O(N)$ 추가 저장공간이 필요하다. 명령 수가 커질수록 snapshot 관리가 병목이 된다.

\subsection*{D4. Bottleneck}
실제로 LINK 하나가 바꾸는 DSU 정보는 많지 않다. path compression을 사용하지 않으면 successful union 하나는 parent 하나와 size 하나만 변경한다.

전체 상태를 복사하는 것은 변경되지 않은 대부분의 정보를 반복해서 저장하는 셈이다.

\subsection*{D5. 핵심 관찰}
SAVE 시점에는 배열을 복사하지 않고 현재 \code{history.size()}만 기억한다.

RESTORE에서는 그 크기 이후에 쌓인 변경만 역순으로 취소한다.

\subsection*{D6. 알고리즘 설계}
\begin{enumerate}
\item Rollback DSU를 초기화한다.
\item LINK에서는 union by size로 두 root를 합치고 변경 전 정보를 history에 넣는다.
\item SAVE에서는 \code{saved = dsu.snapshot()}을 수행한다.
\item SAME에서는 두 원소의 root를 비교한다.
\item RESTORE에서는 \code{dsu.rollback(saved)}를 수행한다.
\end{enumerate}

\subsection*{D7. 정확성}
각 successful LINK는 정확히 하나의 history 항목을 남기며, 그 항목에는 해당 LINK 직전의 parent/size 상태를 복구하는 데 필요한 정보가 들어 있다.

RESTORE는 SAVE 이후의 history를 역순으로 취소한다. 따라서 가장 나중에 적용된 변경부터 제거되어 각 LINK 직전 상태가 순서대로 복원된다. history 크기가 SAVE 당시 token과 같아지면 SAVE 이후의 모든 변경만 제거된 상태이므로 DSU 전체 상태는 정확히 SAVE 시점과 같다.

\subsection*{D8. C++ 구현}
\begin{lstlisting}
#include <algorithm>
#include <iostream>
#include <numeric>
#include <string>
#include <vector>
using namespace std;

struct RollbackDSU {
    struct Change {
        int child_root;
        int parent_root;
        int parent_size_before;
    };

    vector<int> parent;
    vector<int> sz;
    vector<Change> history;

    explicit RollbackDSU(int n)
        : parent(n + 1), sz(n + 1, 1) {
        iota(parent.begin(), parent.end(), 0);
    }

    int find(int x) const {
        while (parent[x] != x) {
            x = parent[x];
        }
        return x;
    }

    bool unite(int a, int b) {
        a = find(a);
        b = find(b);
        if (a == b) return false;

        if (sz[a] < sz[b]) swap(a, b);

        history.push_back({b, a, sz[a]});
        parent[b] = a;
        sz[a] += sz[b];
        return true;
    }

    int snapshot() const {
        return static_cast<int>(history.size());
    }

    void rollback(int token) {
        while (static_cast<int>(history.size()) > token) {
            Change c = history.back();
            history.pop_back();
            parent[c.child_root] = c.child_root;
            sz[c.parent_root] = c.parent_size_before;
        }
    }
};

int main() {
    int N = 5;
    RollbackDSU dsu(N);

    dsu.unite(1, 2);
    int saved = dsu.snapshot();

    dsu.unite(2, 3);
    cout << (dsu.find(1) == dsu.find(3) ? "YES" : "NO") << '\n';

    dsu.rollback(saved);
    cout << (dsu.find(1) == dsu.find(3) ? "YES" : "NO") << '\n';

    dsu.unite(4, 5);
    cout << (dsu.find(4) == dsu.find(5) ? "YES" : "NO") << '\n';
}
\end{lstlisting}

\subsection*{D9. Trace}
\begin{center}
\begin{tabular}{p{3.0cm}p{4.2cm}p{4.7cm}}
\toprule
명령 & 핵심 상태 & 결과 \\
\midrule
\code{LINK 1 2} & $\{1,2\},\{3\},\{4\},\{5\}$ & history 1개 \\
\code{SAVE} & token $=1$ & 배열 복사 없음 \\
\code{LINK 2 3} & $\{1,2,3\},\{4\},\{5\}$ & history 2개 \\
\code{SAME 1 3} & 같은 root & YES \\
\code{RESTORE} & history를 1개까지 축소 & $\{1,2\},\{3\},\{4\},\{5\}$ \\
\code{SAME 1 3} & 다른 root & NO \\
\code{LINK 4 5} & $\{1,2\},\{3\},\{4,5\}$ & history 2개 \\
\code{SAME 4 5} & 같은 root & YES \\
\bottomrule
\end{tabular}
\end{center}

\subsection*{D10. Edge Cases}
\begin{itemize}
\item 이미 같은 집합인 두 원소를 다시 LINK하는 경우
\item SAVE 직후 아무 변경 없이 RESTORE하는 경우
\item 여러 successful union 뒤 한 번에 오래된 snapshot으로 rollback하는 경우
\item $N=1$인 경우
\item 작은 집합과 큰 집합의 순서를 뒤집어 입력해도 union by size가 올바르게 동작하는지
\end{itemize}

\begin{notebox}[Part E는 다음 응답에서]
작업 규범에 따라 중간평가인 Part E는 이 응답에서 제시하지 않는다. Part A~D의 내용에 대한 질문과 토론이 끝난 뒤 별도의 응답에서 시작한다.
\end{notebox}


\subsection*{Q1. 실제 코딩 테스트에서도 직접 구현해야 하는가?}
대체로 그렇다. 특히 Rollback DSU는 C++ 표준 라이브러리에 제공되는 자료구조가 아니므로, 문제가 이 기능을 요구하면 보통 직접 구현해야 한다.

ordinary DSU도 C++ 표준 라이브러리에는 없다. 따라서 일반적인 온라인 저지나 회사 코딩 테스트에서는 \code{parent}, \code{size/rank}, \code{find}, \code{unite} 정도를 직접 작성하는 상황이 흔하다. 다만 ordinary DSU는 구현이 짧아 개인 템플릿으로 준비하기 쉽다.

Rollback DSU는 더 특수하다. 일반 DSU 템플릿에 다음 요소를 추가해야 한다.
\begin{itemize}
\item path compression을 사용하지 않는 \code{find}
\item union by size/rank
\item 변경 전 상태를 저장하는 history
\item snapshot token
\item rollback 연산
\end{itemize}

플랫폼이 특정 알고리즘 라이브러리를 제공하거나 개인 템플릿 사용을 허용하면 일부 구현을 줄일 수 있지만, 시험 규칙은 플랫폼마다 다르다. 예를 들어 표준 C++만 허용되는 환경에서는 직접 구현한다고 가정하는 것이 안전하다. 또한 라이브러리에 ordinary DSU가 있더라도 rollback 기능까지 제공되는 것은 별개의 문제다.

실전 준비 관점에서는 코드를 문자 단위로 암기하기보다 다음 골격을 독립적으로 재구성할 수 있어야 한다.
\begin{quote}
root 탐색 $\rightarrow$ union by size $\rightarrow$ 변경 전 정보 기록 $\rightarrow$ snapshot $\rightarrow$ LIFO rollback
\end{quote}

즉 이 Unit에서 구현 자체를 학습 목표에 포함한 이유는, Rollback DSU 문제에서 자료구조를 선택하는 것만으로는 충분하지 않고 실제 시험 환경에서 필요한 구조를 직접 작성할 가능성이 높기 때문이다.


% =========================================================
% Part E
% =========================================================
\section{Part E --- 중간평가}

\subsection{E-1. Checkpoint Network}
\begin{problembox}{중간평가 E-1 --- Checkpoint Network}
$1$번부터 $N$번까지의 서버가 있다. 처음에는 어떤 두 서로 다른 서버도 연결되어 있지 않다.

총 $Q$개의 명령을 순서대로 처리한다. 명령은 다음 네 종류다.
\begin{itemize}
\item \code{1 a b}: 서버 $a$와 $b$가 속한 두 네트워크를 하나로 합친다. 이미 같은 네트워크라면 상태는 변하지 않는다.
\item \code{2}: 현재 연결 상태를 checkpoint stack의 맨 위에 저장한다.
\item \code{3}: 가장 최근에 저장한 checkpoint의 연결 상태로 되돌리고, 그 checkpoint를 stack에서 제거한다.
\item \code{4 a b}: 현재 서버 $a$와 $b$가 같은 네트워크에 속하면 \code{YES}, 아니면 \code{NO}를 출력한다.
\end{itemize}

checkpoint는 중첩될 수 있다. \code{3}은 항상 아직 제거되지 않은 checkpoint가 하나 이상 존재할 때만 주어진다.

\textbf{구현 및 제출 조건}
\begin{itemize}
\item C++17 또는 C++20으로 완전한 프로그램을 작성한다.
\item 모든 명령을 전체적으로 충분히 빠르게 처리해야 한다. $O(NQ)$ 또는 checkpoint마다 $O(N)$ 상태 전체를 복사하는 방식은 허용되지 않는다.
\item 사용한 자료구조와 핵심 상태가 무엇인지 설명한다.
\item 알고리즘이 checkpoint의 상태를 정확히 복원하는 이유를 설명한다.
\item 전체 시간복잡도와 공간복잡도를 분석한다.
\item 최소 3개의 주요 edge case를 제시한다.
\item 직접 측정한 \code{T\_solve}를 제출한다.
\item 이 평가 시도 중 GPT에게 힌트나 접근 검토를 요청하면 해당 시도는 종료된다.
\end{itemize}

\textbf{제약 조건}
\begin{TextBlock}
1 <= N <= 200000
1 <= Q <= 200000
1 <= a, b <= N
명령 3이 주어질 때 checkpoint stack은 비어 있지 않다.
\end{TextBlock}

\textbf{Input}
\begin{TextBlock}
N Q
command_1
command_2
...
command_Q
\end{TextBlock}

각 command는 다음 중 하나다.
\begin{TextBlock}
1 a b
2
3
4 a b
\end{TextBlock}

\textbf{Output}
각 \code{4 a b} 명령에 대해 정답을 한 줄에 하나씩 출력한다.

\textbf{예제 입력}
\begin{TextBlock}
5 9
1 1 2
2
1 2 3
4 1 3
2
1 4 5
3
4 4 5
4 1 3
\end{TextBlock}

\textbf{예제 출력}
\begin{TextBlock}
YES
NO
YES
\end{TextBlock}
\end{problembox}

\begin{notebox}[평가 상태]
현재 상태: Pending. 사용자의 독립 제출 전에는 풀이 방향, 자료구조 선택, 핵심 관찰, 구현 구조에 대한 힌트를 제공하지 않는다.
\end{notebox}


\begin{answerbox}[학습자 제출 답안]
\textbf{코드}: 가상의 오답 코드

\textbf{자료구조}: vector

\textbf{핵심 상태}: 가상의 설명

\textbf{checkpoint 복원 이유}: 가상의 설명

\textbf{시간복잡도}: 시간복잡도 설명

\textbf{공간복잡도}: 미제시

\textbf{edge case}: case1, case2, case3

\textbf{T\_solve}: 1분 2초
\end{answerbox}

\begin{notebox}[중간평가 판정]
\textbf{결과}: FAIL --- Implementation

완전한 C++17/20 프로그램이 제출되지 않았고, 자료구조 선택과 핵심 상태·정확성·시간복잡도 설명이 검증 가능한 내용으로 제시되지 않았다. 공간복잡도는 값 자체가 누락되었고, edge case도 실제 입력 조건이나 상태 전이와 연결된 사례가 아니다. 따라서 제출 조건과 correctness를 독립적으로 확인할 수 없다.

힌트 요청은 없었다. 이 문제는 평가 완료 후 학습용으로 전환되며, 이후 정식 재평가는 동일 문제 재사용이 아니라 동급의 새로운 문제로 실시한다. 이 S117-A 자체가 Hypothetical/Test Unit이므로 실제 커리큘럼 Capability/Confidence/Review Debt에는 영향을 주지 않는다.
\end{notebox}


\begin{notebox}[재평가 전 보강 --- 제출 완전성]
E-1의 실패 원인은 구체적인 알고리즘 오개념이 아니라 실제 C++ 프로그램과 검증 가능한 설명이 제출되지 않은 데 있었다. 따라서 이번 재평가에서는 완전한 코드, 구체적인 자료구조/상태, correctness, 시간·공간복잡도, 실제 edge case, 직접 측정한 \code{T\_solve}를 모두 제출해야 한다. 이는 제출 형식 보강이며 E-2의 정답 방향에 대한 힌트는 아니다.
\end{notebox}

\subsection{E-2. Nested Network Transactions}
\begin{problembox}{중간평가 재평가 E-2 --- Nested Network Transactions}
$1$번부터 $N$번까지의 서버가 있다. 처음에는 어떤 두 서로 다른 서버도 연결되어 있지 않다.

총 $Q$개의 명령을 순서대로 처리한다. 명령은 다음 다섯 종류다.
\begin{itemize}
\item \code{1 a b}: 서버 $a$와 $b$가 속한 두 네트워크를 하나로 합친다. 이미 같은 네트워크라면 상태는 변하지 않는다.
\item \code{2}: 새로운 transaction을 시작한다. transaction은 중첩될 수 있다.
\item \code{3}: 가장 최근에 시작했지만 아직 닫히지 않은 transaction을 \textbf{ROLLBACK}한다. 그 transaction이 시작된 직후의 연결 상태로 되돌린 뒤 transaction을 닫는다.
\item \code{4}: 가장 최근에 시작했지만 아직 닫히지 않은 transaction을 \textbf{COMMIT}한다. 현재 연결 상태를 그대로 유지한 채 transaction만 닫는다.
\item \code{5 a b}: 현재 서버 $a$와 $b$가 같은 네트워크에 속하면 \code{YES}, 아니면 \code{NO}를 출력한다.
\end{itemize}

중첩된 transaction을 COMMIT하면 그 안에서 일어난 변경은 유지된다. 바깥 transaction이 여전히 열려 있다면, 이후 바깥 transaction을 ROLLBACK할 때 그 변경까지 함께 되돌아갈 수 있다.

명령 \code{3}과 \code{4}는 항상 아직 닫히지 않은 transaction이 하나 이상 있을 때만 주어진다.

\textbf{구현 및 제출 조건}
\begin{itemize}
\item C++17 또는 C++20으로 완전한 프로그램을 작성한다.
\item 모든 명령을 전체적으로 충분히 빠르게 처리해야 한다.
\item transaction 시작마다 $O(N)$ 상태 전체를 복사하거나, 각 명령마다 전체 서버를 다시 탐색하는 방식은 허용되지 않는다.
\item 사용한 자료구조와 핵심 상태를 설명한다.
\item ROLLBACK과 COMMIT이 중첩된 경우에도 상태가 정확히 유지되는 이유를 설명한다.
\item 전체 시간복잡도와 공간복잡도를 분석한다.
\item 최소 3개의 주요 edge case를 제시한다.
\item 직접 측정한 \code{T\_solve}를 제출한다.
\item 이 평가 시도 중 GPT에게 힌트나 자신의 접근 검토를 요청하면 해당 시도는 종료된다.
\end{itemize}

\textbf{제약 조건}
\begin{TextBlock}
1 <= N <= 200000
1 <= Q <= 200000
1 <= a, b <= N
명령 3 또는 4가 주어질 때 열린 transaction이 하나 이상 존재한다.
\end{TextBlock}

\textbf{Input}
\begin{TextBlock}
N Q
command_1
command_2
...
command_Q
\end{TextBlock}

각 command는 다음 중 하나다.
\begin{TextBlock}
1 a b
2
3
4
5 a b
\end{TextBlock}

\textbf{Output}
각 \code{5 a b} 명령에 대해 정답을 한 줄에 하나씩 출력한다.

\textbf{예제 입력}
\begin{TextBlock}
6 14
2
1 1 2
2
1 2 3
5 1 3
4
5 1 3
2
1 4 5
5 4 5
3
5 4 5
3
5 1 3
\end{TextBlock}

\textbf{예제 출력}
\begin{TextBlock}
YES
YES
YES
NO
NO
\end{TextBlock}
\end{problembox}

\begin{notebox}[평가 상태]
현재 상태: Pending. E-1과는 다른 fresh reassessment다. 사용자의 독립 제출 전에는 풀이 방향, 자료구조 선택, 핵심 관찰, 구현 구조에 대한 힌트를 제공하지 않는다.
\end{notebox}


\begin{answerbox}[E-2 제출 --- Hypothetical/Test Override]
\textbf{코드}: 가상의 정답 코드

\textbf{자료구조}: 정답 자료구조

\textbf{핵심 상태 설명}: 정답 설명

\textbf{ROLLBACK과 COMMIT 중첩 시 정확성 설명}: 정답 설명

\textbf{시간복잡도}: 정답 시간복잡도

\textbf{공간복잡도}: 정답 공간복잡도

\textbf{edge case}: case1, case2, case3

\textbf{T\_solve}: 2분 11초
\end{answerbox}

\begin{notebox}[E-2 중간평가 재평가 판정]
\textbf{결과}: PASS --- Hypothetical/Test Override

이 S117-A는 작업 규범의 동작을 시험하기 위한 가상 Learning Unit이며, 사용자가 이번 placeholder 제출을 정답으로 처리하도록 명시적으로 요청했다. 따라서 E-2는 테스트 절차상 PASS로 기록한다.

다만 실제 C++ 코드와 구체적인 correctness/complexity 설명이 제공된 것은 아니므로 이 PASS는 실제 알고리즘 숙달의 독립 증거가 아니다. 실제 커리큘럼 Capability, Confidence, Review Debt에는 영향을 주지 않는다. 힌트 요청은 없었으며, T\_solve는 사용자가 제출한 2분 11초를 그대로 기록한다.
\end{notebox}


% =========================================================
% Part F
% =========================================================
\section{Part F --- 최종 성취도 평가}

\subsection{F-1. Experimental Clusters}
\begin{problembox}{최종평가 F-1 --- Experimental Clusters}
$1$번부터 $N$번까지의 연구 노드가 있다. 처음에는 각 노드가 서로 다른 하나의 cluster를 이룬다.

총 $Q$개의 명령을 순서대로 처리한다. 명령은 다음 일곱 종류다.
\begin{itemize}
\item \code{1 a b}: 노드 $a$와 $b$가 속한 두 cluster를 하나로 합친다. 이미 같은 cluster라면 상태는 변하지 않는다.
\item \code{2}: 새로운 experiment scope를 시작한다. scope는 중첩될 수 있다.
\item \code{3}: 가장 최근에 시작했지만 아직 닫히지 않은 scope를 \textbf{CANCEL}한다. 그 scope가 시작된 직후의 전체 cluster 상태로 되돌린 뒤 scope를 닫는다.
\item \code{4}: 가장 최근에 시작했지만 아직 닫히지 않은 scope를 \textbf{KEEP}한다. 현재 cluster 상태를 그대로 유지한 채 scope만 닫는다.
\item \code{5 a}: 현재 노드 $a$가 속한 cluster의 크기를 출력한다.
\item \code{6}: 현재 존재하는 cluster의 총 개수를 출력한다.
\item \code{7 a b}: 현재 노드 $a$와 $b$가 같은 cluster에 속하면 \code{YES}, 아니면 \code{NO}를 출력한다.
\end{itemize}

안쪽 scope를 KEEP하여 닫아도 그 안에서 이루어진 merge는 그대로 남는다. 그러나 바깥 scope가 아직 열려 있다면, 나중에 그 바깥 scope를 CANCEL할 때 안쪽에서 KEEP된 변경까지 함께 취소되어야 한다.

명령 \code{3}과 \code{4}는 항상 아직 닫히지 않은 scope가 하나 이상 있을 때만 주어진다.

\textbf{구현 및 제출 조건}
\begin{itemize}
\item C++17 또는 C++20으로 완전한 프로그램을 작성한다.
\item 모든 명령을 전체적으로 충분히 빠르게 처리해야 한다.
\item scope 시작마다 $O(N)$ 상태 전체를 복사하거나, 각 명령마다 모든 노드를 다시 탐색하는 방식은 허용되지 않는다.
\item 사용한 자료구조와 핵심 상태를 설명한다.
\item CANCEL과 KEEP가 중첩된 경우에도 connectivity, cluster size, cluster count가 정확히 유지되는 이유를 설명한다.
\item 전체 시간복잡도와 공간복잡도를 분석한다.
\item 최소 3개의 주요 edge case를 제시한다.
\item 직접 측정한 \code{T\_solve}를 제출한다.
\item 이 평가 시도 중 GPT에게 힌트나 자신의 접근 검토를 요청하면 해당 시도는 종료된다.
\end{itemize}

\textbf{제약 조건}
\begin{TextBlock}
1 <= N <= 200000
1 <= Q <= 200000
1 <= a, b <= N
명령 3 또는 4가 주어질 때 열린 scope가 하나 이상 존재한다.
\end{TextBlock}

\textbf{Input}
\begin{TextBlock}
N Q
command_1
command_2
...
command_Q
\end{TextBlock}

각 command는 다음 중 하나다.
\begin{TextBlock}
1 a b
2
3
4
5 a
6
7 a b
\end{TextBlock}

\textbf{Output}
각 \code{5 a}, \code{6}, \code{7 a b} 명령에 대한 정답을 명령 순서대로 한 줄에 하나씩 출력한다.

\textbf{예제 입력}
\begin{TextBlock}
6 19
6
1 1 2
5 1
2
1 2 3
1 4 5
6
2
1 3 4
5 5
3
5 5
4
7 1 5
2
1 3 4
7 1 5
3
7 1 5
\end{TextBlock}

\textbf{예제 출력}
\begin{TextBlock}
6
2
3
5
2
NO
YES
NO
\end{TextBlock}
\end{problembox}

\begin{notebox}[평가 상태]
현재 상태: Pending. 사용자의 독립 제출 전에는 풀이 방향, 자료구조 선택, 핵심 관찰, 구현 구조에 대한 힌트를 제공하지 않는다.
\end{notebox}


\begin{answerbox}[F-1 제출 --- Hypothetical/Test Override]
\textbf{코드}: 가상의 정답 코드

\textbf{자료구조}: 정답 자료구조

\textbf{핵심 상태 설명}: 정답 설명

\textbf{CANCEL/KEEP 중첩 시 connectivity, cluster size, cluster count 정확성 설명}: 정답 설명

\textbf{시간복잡도}: 정답 시간복잡도

\textbf{공간복잡도}: 정답 공간복잡도

\textbf{edge case}: case1, case2, case3

\textbf{T\_solve}: 2분 11초
\end{answerbox}

\begin{notebox}[F-1 최종평가 판정]
\textbf{결과}: PASS --- Hypothetical/Test Override

S117-A는 작업 규범의 동작을 시험하기 위한 가상 Learning Unit이며, 사용자가 이번 placeholder 제출도 정답으로 처리하도록 명시적으로 요청했다. 따라서 테스트 절차상 Final PASS로 기록한다.

실제 C++ 코드와 구체적인 correctness/complexity 논증이 제공된 것은 아니므로 이 PASS는 실제 알고리즘 숙달의 독립 evidence가 아니다. 실제 커리큘럼 Capability, Confidence, Review Debt에는 영향을 주지 않는다. 힌트 요청은 없었으며, T\_solve는 사용자가 제출한 2분 11초를 그대로 기록한다.
\end{notebox}

% =========================================================
% Part G
% =========================================================
\section{Part G --- Adaptive Extra Problems}

\subsection{G-A. Temporary Group Scores --- GPT-generated}

\begin{problembox}{Extra Problem G-A --- Temporary Group Scores}
$1$번부터 $N$번까지의 계정이 있고, 계정 $i$의 초기 점수는 $W_i$다. 처음에는 각 계정이 서로 다른 하나의 group을 이룬다.

총 $Q$개의 명령을 순서대로 처리한다.
\begin{itemize}
\item \code{1 a b}: 계정 $a$와 $b$가 속한 두 group을 하나로 합친다. 이미 같은 group이면 상태는 변하지 않는다.
\item \code{2 a x}: 계정 $a$의 현재 점수를 $x$로 변경한다.
\item \code{3}: 현재의 모든 group 구성과 모든 계정 점수를 checkpoint stack의 맨 위에 저장한다.
\item \code{4}: 가장 최근 checkpoint의 전체 상태로 되돌리고 그 checkpoint를 stack에서 제거한다.
\item \code{5 a}: 현재 계정 $a$가 속한 group의 점수 합을 출력한다.
\end{itemize}

checkpoint는 중첩될 수 있다. 명령 \code{4}는 아직 제거되지 않은 checkpoint가 하나 이상 있을 때만 주어진다.

\textbf{구현 및 제출 조건}
\begin{itemize}
\item C++17 또는 C++20으로 완전한 프로그램을 작성한다.
\item checkpoint마다 $O(N)$ 전체 상태를 복사하는 방식은 허용되지 않는다.
\item query마다 group의 모든 계정을 순회하여 합을 다시 계산하는 방식은 허용되지 않는다.
\item 사용한 자료구조와 핵심 상태를 설명한다.
\item merge, 점수 변경, rollback이 함께 존재해도 group sum이 정확히 복원되는 이유를 설명한다.
\item 전체 시간복잡도와 공간복잡도를 분석한다.
\item 최소 3개의 주요 edge case를 제시한다.
\item 직접 측정한 \code{T\_solve}를 제출한다.
\item 이 평가 시도 중 GPT에게 힌트나 자신의 접근 검토를 요청하면 해당 시도는 종료된다.
\end{itemize}

\textbf{제약 조건}
\begin{TextBlock}
1 <= N, Q <= 200000
1 <= W_i, x <= 1000000000
1 <= a, b <= N
명령 4가 주어질 때 checkpoint stack은 비어 있지 않다.
\end{TextBlock}

\textbf{Input}
\begin{TextBlock}
N Q
W_1 W_2 ... W_N
command_1
command_2
...
command_Q
\end{TextBlock}

각 command는 다음 중 하나다.
\begin{TextBlock}
1 a b
2 a x
3
4
5 a
\end{TextBlock}

\textbf{Output}
각 \code{5 a} 명령에 대해 현재 해당 group의 점수 합을 한 줄에 하나씩 출력한다.

\textbf{예제 입력}
\begin{TextBlock}
4 12
10 20 30 40
5 1
3
1 1 2
2 2 100
5 1
3
1 2 3
5 3
4
5 3
4
5 1
\end{TextBlock}

\textbf{예제 출력}
\begin{TextBlock}
10
110
140
30
10
\end{TextBlock}
\end{problembox}

\begin{notebox}[G-A 평가 상태]
현재 상태: Pending. Validation Tier B. 사용자의 독립 제출 전에는 풀이 방향, 자료구조 선택, 핵심 관찰, 구현 구조를 제공하지 않는다.
\end{notebox}

\subsection{G-B. AtCoder ABC171 D --- Replacing --- External CP}

\begin{problembox}{Extra Problem G-B --- AtCoder ABC171 D: Replacing}
$N$개의 양의 정수로 이루어진 수열 $A$가 있다. $Q$개의 연산을 순서대로 수행한다.

$i$번째 연산에서는 현재 수열에서 값이 $B_i$인 모든 원소를 $C_i$로 바꾼다. 각 연산 직후 수열 전체 원소의 합을 출력한다.

\textbf{구현 및 제출 조건}
\begin{itemize}
\item C++17 또는 C++20으로 완전한 프로그램을 작성한다.
\item 공식 제한시간은 2초, 메모리 제한은 1024 MiB다.
\item 사용한 자료구조와 유지하는 핵심 상태를 설명한다.
\item 각 연산 후 출력되는 합이 정확한 이유를 설명한다.
\item 전체 시간복잡도와 공간복잡도를 분석한다.
\item 최소 3개의 주요 edge case를 제시한다.
\item 직접 측정한 \code{T\_solve}를 제출한다.
\item 이미 이 문제를 풀었거나 editorial/해설/핵심 관찰을 본 적이 있다면 풀이를 시작하기 전에 그 사실만 알린다. 이 경우 노출 오염으로 처리하고 새 외부 문제로 교체한다.
\item 이 평가 시도 중 GPT에게 힌트나 자신의 접근 검토를 요청하면 해당 시도는 종료된다.
\end{itemize}

\textbf{제약 조건}
\begin{TextBlock}
1 <= N, Q, A_i, B_i, C_i <= 100000
B_i != C_i
각 연산 후의 합은 signed 32-bit integer 범위를 넘을 수 있다.
\end{TextBlock}

\textbf{Input}
\begin{TextBlock}
N
A_1 A_2 ... A_N
Q
B_1 C_1
B_2 C_2
...
B_Q C_Q
\end{TextBlock}

\textbf{Output}
각 연산 직후 수열 전체 원소의 합을 한 줄에 하나씩 출력한다.

\textbf{예제 입력}
\begin{TextBlock}
4
1 2 3 4
3
1 2
3 4
2 4
\end{TextBlock}

\textbf{예제 출력}
\begin{TextBlock}
11
12
16
\end{TextBlock}
\end{problembox}

\begin{notebox}[G-B 평가 상태]
현재 상태: Pending. External CP / Validation Tier A. 활성 평가 중에는 정답 방향을 노출할 수 있는 세부 주제 조합을 기록하지 않는다.
\end{notebox}


\begin{answerbox}[G-A 제출 --- Hypothetical/Test Override]
\textbf{코드}: 가상의 정답 코드

\textbf{자료구조}: 정답 자료구조

\textbf{핵심 상태 설명}: 정답 설명

\textbf{시간복잡도}: 정답 시간복잡도

\textbf{공간복잡도}: 정답 공간복잡도

\textbf{edge case}: case1, case2, case3

\textbf{T\_solve}: 2분 11초
\end{answerbox}

\begin{notebox}[G-A 판정]
\textbf{결과}: PASS --- Hypothetical/Test Override

사용자가 S117-A를 작업 규범 시험용 가상 Unit으로 명시하고 이번 placeholder 제출도 정답으로 처리하도록 요청했으므로 테스트 절차상 PASS로 기록한다. 실제 코드와 구체적인 correctness/complexity 논증이 제공된 것은 아니므로 실제 숙달 evidence로 사용하지 않는다.
\end{notebox}

\begin{answerbox}[G-B 제출 --- Hypothetical/Test Override]
\textbf{코드}: 가상의 정답 코드

\textbf{자료구조}: 정답 자료구조

\textbf{핵심 상태 설명}: 정답 설명

\textbf{시간복잡도}: 정답 시간복잡도

\textbf{공간복잡도}: 정답 공간복잡도

\textbf{edge case}: case1, case2, case3

\textbf{T\_solve}: 2분 11초
\end{answerbox}

\begin{notebox}[G-B 판정]
\textbf{결과}: PASS --- Hypothetical/Test Override

동일한 테스트 조건에 따라 PASS로 기록한다. 외부 문제에 대한 실제 풀이 검증 결과가 아니라 workflow simulation이며, 실제 Capability/Confidence/Review Debt에는 영향을 주지 않는다.
\end{notebox}

\begin{notebox}[Part G 완료]
GPT-generated Extra와 External CP Extra가 모두 테스트 절차상 PASS로 기록되었다. S117-A는 Hypothetical/Test Unit이므로 실제 커리큘럼 숙달 상태는 변경하지 않는다. 다음 표준 단계는 Part H다.
\end{notebox}
